\documentclass{umk-matematika}
\begin{document}

\begin{center}
{\Large\bfseries Практическая работа 28}
\end{center}
\vspace{0.5cm}

\textbf{Задача 1.} Радиус основания конуса $6$ см, образующая $10$ см. Найдите площадь боковой поверхности.\par\vspace{10pt}
\textbf{Задача 2.} Радиус основания конуса $4$ см, высота $3$ см. Найдите объём.\par\vspace{10pt}
\textbf{Задача 3.} Площадь боковой поверхности конуса $45\pi$ см², радиус основания $5$ см. Найдите образующую.\par\vspace{10pt}
\textbf{Задача 4.} Объём конуса $100\pi$ см³, радиус основания $5$ см. Найдите высоту.\par\vspace{10pt}
\textbf{Задача 5.} Осевое сечение конуса — равносторонний треугольник со стороной $8$ см. Найдите площадь полной поверхности конуса.\par\vspace{10pt}
\textbf{Задача 6.} Образующая конуса $13$ см, высота $12$ см. Найдите объём конуса.\par\vspace{10pt}
\textbf{Задача 7.} Радиус основания конуса $7$ см, высота $24$ см. Найдите площадь полной поверхности.\par\vspace{10pt}
\textbf{Задача 8.} Куча песка имеет форму конуса. Длина окружности основания $6\pi$ м, образующая $5$ м. Найдите объём песка.\par\vspace{10pt}
\textbf{Задача 9.} Крыша башни — конус с диаметром основания $10$ м и высотой $12$ м. Сколько листов кровли нужно, если один лист покрывает $2$ м²? (С запасом $10\\%$)\par\vspace{10pt}
\textbf{Задача 10.} Резервуар состоит из цилиндра (радиус $3$ м, высота $4$ м) и конуса (такой же радиус, высота $2$ м), поставленного сверху. Найдите общий объём.\par\vspace{10pt}
\newpage
\section*{Ответы}

\begin{enumerate}
  \item Ответ: $60\pi$ см²
  \item Ответ: $16\pi$ см³
  \item Ответ: $9$ см
  \item Ответ: $12$ см
  \item Ответ: $48\pi$ см²
  \item Ответ: $100\pi$ см³
  \item Ответ: $224\pi$ см²
  \item Ответ: $12\pi$ м³ $\approx 37{,}7$ м³
  \item Ответ: $113$ листов
  \item Ответ: $42\pi$ м³ $\approx 131{,}9$ м³
\end{enumerate}
\end{document}
